\documentclass[11pt]{article}
\usepackage[T1]{fontenc}
\usepackage{lmodern}
\usepackage[margin=1in]{geometry}
\usepackage{amsmath,amssymb,amsthm,mathtools}
\usepackage{microtype,booktabs,longtable,tabularx,array}
\usepackage{enumitem}
\usepackage{xcolor}
\usepackage[hyphens]{url}
\usepackage[colorlinks=true,linkcolor=black,citecolor=black,urlcolor=black]{hyperref}
\usepackage{fancyhdr}
\setlength{\headheight}{14pt}
\pagestyle{fancy}
\fancyhf{}
\fancyhead[L]{\small Signed-kernel oscillation at every depth}
\fancyhead[R]{\small ProveIt research continuation}
\fancyfoot[C]{\thepage}
\setlength{\parskip}{4pt}
\setlength{\parindent}{1.2em}
\setlist{topsep=3pt,itemsep=3pt,parsep=0pt}
\emergencystretch=2em
\numberwithin{equation}{section}
\newtheorem{theorem}{Theorem}[section]
\newtheorem{lemma}[theorem]{Lemma}
\newtheorem{proposition}[theorem]{Proposition}
\newtheorem{corollary}[theorem]{Corollary}
\newtheorem{conjecture}[theorem]{Conjecture}
\theoremstyle{definition}
\newtheorem{definition}[theorem]{Definition}
\newtheorem{remark}[theorem]{Remark}
\DeclareMathOperator{\Li}{Li}
\DeclareMathOperator{\Imn}{Im}
\DeclareMathOperator{\Ren}{Re}
\DeclareMathOperator{\Arg}{Arg}
\DeclareMathOperator{\supp}{supp}
\newcommand{\dd}{\,\mathrm d}
\newcommand{\ii}{\mathrm i}
\newcommand{\bbR}{\mathbb R}
\newcommand{\bbC}{\mathbb C}
\newcommand{\bbN}{\mathbb N}
\newcommand{\svec}{\boldsymbol{s}}
\newcommand{\one}{\boldsymbol{1}}
\newcommand{\Om}{\Omega}
\newcommand{\Aop}{\mathsf A}
\newcommand{\Bop}{\mathsf B}
\newcommand{\norm}[1]{\lVert#1\rVert}
\newcommand{\abs}[1]{\lvert#1\rvert}
\newcommand{\bracks}[2]{\genfrac{[}{]}{0pt}{}{#1}{#2}}
\newcommand{\repoSHA}{3447bc59b78d138a7f0d536b66ac6e5cc5d5e49f}
\hypersetup{pdftitle={Signed-Kernel Oscillation at Every Depth},pdfauthor={OpenAI-assisted research continuation for ProveIt},pdfsubject={Strict multiple polylogarithms, positive compensation, angular zeros, exact certificates}}
\title{\textbf{Signed-Kernel Oscillation at Every Depth}\\[6pt]
\large Minimal Pick Compensation, Zero Geometry,\\and Certified Polylogarithmic Identities}
\author{Research continuation prepared for the ProveIt project\\
\small OpenAI-assisted mathematical research note}
\date{October 10, 2026}
\begin{document}
\maketitle
\thispagestyle{empty}
\begin{abstract}
We extend the manuscript's depth-two signed-kernel method to strict
single-variable multiple polylogarithms of arbitrary depth. The first
$d-1$ indices are positive integers and the final index may be any positive
real number. Two bounded operators on $L^1(0,1)$ construct a density with
exactly $d-1$ simple sign changes, $d-1$ vanishing initial moments, and
norm at most $2^{d-1}$. The sign-change polynomial converts this signed
measure into a positive measure and yields a canonical factorization
$L_{\svec}(z)=z^dH_{\svec}(z)/\prod_j(1-r_jz)$. We prove that this
polynomial compensation is of minimal degree, that the principal
continuation has no nonzero zero on $\bbC\setminus[1,\infty)$, and that
its imaginary part has exactly $d-1$ simple zeros on every upper
semicircle of radius at most one. A direct confluent Cauchy-kernel argument
proves the upper zero bound; a positive-transform argument supplies the
matching lower bound, including the singular unit-radius endpoint.

Further results include strict inward angular motion, interlacing of density
zeros, small-radius and
large-outer-index expansions, positive terminal mixtures, a shifted-terminal
Lerch extension, explicit logistic-root identities for the all-ones family,
and rational error certificates at Gaussian points. The package contains
24 certified angular-root brackets and 5,713 finite exact check groups.
The proofs are ordinary analytic proofs, not proof-assistant formalizations.
The $S_6$ and $S_8$ arithmetic conjectures, unrestricted fractional leading
indices, and the existing normalized-radius conjecture are not claimed
solved. No universal claim of literature priority is made.
\end{abstract}
\vfill
\noindent\small\textbf{Repository reference.}
\url{https://github.com/VladimirReshetnikov/ProveIt}\\
Reviewed snapshot: \texttt{\repoSHA}. This is an additive research package;
it does not alter the repository or replace its canonical manuscript.
\clearpage
\setcounter{tocdepth}{1}
\hypersetup{bookmarksdepth=2}
\tableofcontents
\vspace{1.5em}
\noindent\textbf{Reading map.} Sections 2--5 contain the complete structural proof.
Section 6 derives asymptotic consequences, Section 7 gives explicit identities,
and Section 8 explains the exact certificates. Sections 9--10 separate
corrections, scope limits, and future research. The appendices record
independent arithmetic controls and integration dependencies.
\clearpage
\section{The problem, conventions, and contribution}
\label{ado:sec:scope}
The existing ProveIt manuscript develops a successful bridge between
nested sums, signed Hausdorff moments, positive transforms, and angular
zero geometry. Its depth-two function is
\[
 F_{a,b}(z)=\sum_{n\ge1}\frac{H_{n-1}^{(b)}}{n^a}z^n,
 \qquad H_{n-1}^{(b)}=\sum_{m=1}^{n-1}m^{-b}.
\]
In the reviewed snapshot, slit-plane nonvanishing and upper-semicircle
angular uniqueness are already proved at every positive inner order and
nonnegative outer order. The integer-order signed density and the later
real-order and subcritical extensions are distinct parts of that
argument~\cite{repoSigned,repoReal,repoSubcritical}. We do not present
these depth-two conclusions as newly resolved conjectures.

The question pursued here is what survives at arbitrary \emph{strict}
depth. A positive integral with a product of several resolvent factors
is useful but does not, by itself, imply nonvanishing. The missing
structure is a signed density whose number of sign changes matches
exactly the number of initially vanishing coefficients. This equality
allows polynomial compensation and controls the full angular zero count.

\subsection{Strict single-variable notation}
For an index of depth $d$, define
\begin{equation}\label{ado:eq:series}
 L_{\svec}(z)=\Li_{s_1,\ldots,s_d}(z,1,\ldots,1)
 =\sum_{n_1>\cdots>n_d\ge1}
       \frac{z^{n_1}}{n_1^{s_1}\cdots n_d^{s_d}}
 =\sum_{n\ge1}c_{\svec}(n)z^n,\qquad |z|<1.
\end{equation}
Our parameter domain is
\begin{equation}\label{ado:eq:domain}
 s_1,\ldots,s_{d-1}\in\bbN_{\ge1},\qquad s_d=b>0.
\end{equation}
At depth one, this means any $b>0$. Indices are strictly decreasing, not
weakly decreasing. Only the outer summation variable receives $z^{n_1}$.
This convention agrees with the single-variable convention in the
manuscript and in the multiple-polylogarithm literature. Classical
polylogarithm branch conventions are recorded in DLMF~\cite{DLMF}.

Write
\begin{equation}\label{ado:eq:leading}
 C_{\svec}=c_{\svec}(d)
 =\frac{1}{d^{s_1}(d-1)^{s_2}\cdots 2^{s_{d-1}}}>0.
\end{equation}
The empty product at depth one is one. The value of $C_{\svec}$ is
independent of $b$, because the innermost index of the first admissible
chain is one. Also $c_{\svec}(1)=\cdots=c_{\svec}(d-1)=0$.
Throughout, principal continuation means continuation to
\[
 \Om=\bbC\setminus[1,\infty).
\]
The word ``zero'' refers to a zero of a holomorphic function, unless the
imaginary part or a signed density is explicitly specified.

\subsection{Main result}
\begin{theorem}[All-depth compensation and zero geometry]
\label{ado:thm:main}
Let \eqref{ado:eq:domain} hold. There is a smooth real density
$\kappa_{\svec}\in L^1(0,1)$ such that
\[
 \int_0^1u^{n-1}\kappa_{\svec}(u)\dd u=c_{\svec}(n),\qquad
 \norm{\kappa_{\svec}}_1\le2^{d-1}.
\]
It has exactly $d-1$ simple zeros $0<r_1<\cdots<r_{d-1}<1$, with
alternating signs and positive sign near one. Set
\[
 P_{\svec}(u)=\prod_{j=1}^{d-1}(u-r_j),\qquad
 D_{\svec}(z)=\prod_{j=1}^{d-1}(1-r_jz).
\]
Then $\dd\mu_{\svec}=P_{\svec}\kappa_{\svec}\dd u$ is a positive
measure of mass $C_{\svec}$, and
\begin{equation}\label{ado:eq:mainfactor}
 L_{\svec}(z)=\frac{z^d}{D_{\svec}(z)}H_{\svec}(z),\qquad
 H_{\svec}(z)=\int_0^1\frac{\dd\mu_{\svec}(u)}{1-zu}.
\end{equation}
The polynomial $D_{\svec}$ is the unique normalized polynomial of
smallest degree making $D_{\svec}L_{\svec}/z^d$ a positive
Hausdorff--Stieltjes transform. In particular, $L_{\svec}$ has no zero
in $\Om$ other than its zero of multiplicity $d$ at zero.

For every $0<\rho\le1$, there are exactly $d-1$ simple zeros
\[
 0<\theta_1(\rho)<\cdots<\theta_{d-1}(\rho)<\pi
\]
of $\theta\mapsto\Imn L_{\svec}(\rho e^{\ii\theta})$.
With $\alpha_k=k\pi/d$ they satisfy
\begin{equation}\label{ado:eq:mainangles}
 \alpha_k-\arcsin(\rho\sin\alpha_k)
 <\theta_k(\rho)<\alpha_k.
\end{equation}
Every branch satisfies $\theta_k'(\rho)<0$ for $0<\rho\le1$ (with the
one-sided derivative at one). The imaginary part starts positive and
alternates sign at these zeros;
the value of $L_{\svec}$ at the $k$th crossing is real with sign $(-1)^k$.
\end{theorem}
The proof is divided into operator, oscillation, compensation, and
angular-count steps below. None of those steps is inferred from the
finite verification suite.

\subsection{Relation to prior work and what is new here}
Positive Hausdorff generating functions have a classical Pick-function
interpretation; Liu and Pego give a systematic treatment~\cite{LiuPego}.
The present $H_{\svec}$ belongs to that class. The substantial step is
not that a positive measure has a Pick transform, but the construction
of the \emph{minimal} polynomial that changes a strict-depth signed
measure into a positive one.

Recent work on multiple \emph{star} polylogarithms also uses positive
Hausdorff measures~\cite{LiStar}. A star sum has weak inequalities in
its nested indices. It must not be substituted for
\eqref{ado:eq:series}: its first nonzero coefficient, normalization,
and positive-measure properties are different. Our proof treats strict
sums directly and retains their $d-1$ vanishing initial moments.

Relative to the reviewed manuscript sections, the added results are the
two-operator all-depth construction, exact oscillation and interlacing,
minimal positive compensation, the all-depth slit-plane and angular-zero
theorems, their uniform asymptotic consequences, and explicit
compensated integral identities. The elementary all-ones logarithm
formula and the general positive-transform theory are background, not
claims of discovery. A targeted literature search did not establish
priority for every theorem stated here; a broader novelty review remains
appropriate before publication.

\begin{center}\small
\begin{tabularx}{\linewidth}{@{}p{0.30\linewidth}X@{}}
\toprule
Claim & Status and scope\\
\midrule
All-depth kernel and zero count & Proved here on \eqref{ado:eq:domain}.\\
Canonical compensator & Proved minimal and unique among normalized least-degree polynomials.\\
Explicit integral identities & Proved; the underlying all-ones polylogarithm is classical.\\
Gaussian and root certificates & Finite exact rational checks with analytic error bounds.\\
Arbitrary fractional leading orders & Not covered by this proof.\\
$S_6$, new $S_8$, period independence & Not settled or promoted by this package.\\
\bottomrule
\end{tabularx}
\end{center}

\section{Two bounded operators on signed moment densities}
\label{ado:sec:operators}
For $f\in L^1(0,1)$ define its moments by
\[
 m_n(f)=\int_0^1u^{n-1}f(u)\dd u,\qquad n\ge1.
\]
We use the two operators
\begin{align}
 (\Aop f)(u)&=\int_u^1\frac{f(v)}v\dd v,
 \label{ado:eq:A}\\
 (\Bop f)(u)&=\int_0^u\frac{f(v)}{1-v}\dd v
              -\int_u^1\frac{f(v)}v\dd v.
 \label{ado:eq:B}
\end{align}
For every fixed $0<u<1$, both integrals in \eqref{ado:eq:B} are
absolutely convergent. In particular, this definition never subtracts
two divergent endpoint integrals. The operator $\Aop$ is the
outer-index-raising operator already present at depth two in the
manuscript~\cite{repoSigned}. The companion $\Bop$ encodes a strict
prefix sum of moments.

\begin{lemma}[Norm, derivative, and moment identities]
\label{ado:lem:operators}
Both operators map $L^1(0,1)$ to itself and satisfy
\begin{equation}\label{ado:eq:opnorm}
 \norm{\Aop f}_1\le\norm f_1,\qquad
 \norm{\Bop f}_1\le2\norm f_1.
\end{equation}
Their moments are
\begin{equation}\label{ado:eq:opmoments}
 m_n(\Aop f)=\frac{m_n(f)}n,\qquad
 m_n(\Bop f)=\frac1n\sum_{j=1}^{n-1}m_j(f).
\end{equation}
If $f$ is continuous in the open interval, then
\begin{equation}\label{ado:eq:opderivatives}
 (\Aop f)'(u)=-\frac{f(u)}u,\qquad
 (\Bop f)'(u)=\frac{f(u)}{u(1-u)}.
\end{equation}
Moreover $\Aop f$ extends continuously to one with value zero, and
$m_1(\Bop f)=0$.
\end{lemma}
\begin{proof}
Tonelli's theorem applied to absolute values gives
\[
 \int_0^1\int_u^1\frac{|f(v)|}{v}\dd v\dd u
 =\int_0^1|f(v)|\dd v.
\]
Similarly,
\[
 \int_0^1\int_0^u\frac{|f(v)|}{1-v}\dd v\dd u
 =\int_0^1|f(v)|\dd v.
\]
These identities prove the two operator bounds and justify every
following application of Fubini to a signed integrand. For the first
operator,
\[
 m_n(\Aop f)=\int_0^1\frac{f(v)}v
       \left(\int_0^v u^{n-1}\dd u\right)\dd v
 =\frac{m_n(f)}n.
\]
The moment of the first term of $\Bop f$ is
\[
 \frac1n\int_0^1\frac{1-v^n}{1-v}f(v)\dd v
 =\frac1n\sum_{j=1}^n m_j(f).
\]
Subtracting $m_n(f)/n$ yields the strict prefix in
\eqref{ado:eq:opmoments}. The derivative formulas follow locally from
the fundamental theorem of calculus. Finally,
$\int_u^1f(v)/v\dd v\to0$ as $u\uparrow1$, because $1/v$ is bounded
near one and $f$ is integrable there.
\end{proof}

\subsection{Construction at arbitrary depth}
The positive starting density is
\begin{equation}\label{ado:eq:seed}
 f_b(u)=\frac{(-\log u)^{b-1}}{\Gamma(b)},\qquad b>0.
\end{equation}
Its mass is one and the Gamma integral gives $m_n(f_b)=n^{-b}$.
Define, with composition read from right to left,
\begin{equation}\label{ado:eq:kappaconstruction}
 \kappa_{\svec}
 =\Aop^{s_1-1}\Bop\Aop^{s_2-1}\Bop\cdots
       \Aop^{s_{d-1}-1}\Bop f_b.
\end{equation}
There are $d-1$ occurrences of $\Bop$; all powers of $\Aop$ are
nonnegative integer powers. At depth one the right side is simply $f_b$.

\begin{proposition}[Moment realization and continuation]
\label{ado:prop:realization}
The density \eqref{ado:eq:kappaconstruction} is smooth on $(0,1)$,
belongs to $L^1(0,1)$, and satisfies
\begin{equation}\label{ado:eq:moments}
 m_n(\kappa_{\svec})=c_{\svec}(n),\qquad
 \norm{\kappa_{\svec}}_1\le2^{d-1}.
\end{equation}
The principal continuation is
\begin{equation}\label{ado:eq:signedtransform}
 L_{\svec}(z)=z\int_0^1\frac{\kappa_{\svec}(u)}{1-zu}\dd u,
 \qquad z\in\Om.
\end{equation}
The density is the unique finite signed measure realizing these moments.
\end{proposition}
\begin{proof}
The coefficient recurrence for a nonempty tail $\boldsymbol{t}$ is
\[
 c_{(a,\boldsymbol{t})}(n)
 =\frac1{n^a}\sum_{j=1}^{n-1}c_{\boldsymbol{t}}(j).
\]
It is exactly the effect of $\Aop^{a-1}\Bop$ in
\eqref{ado:eq:opmoments}. Induction proves the moments and the norm
bound. The derivative identities prove local smoothness, starting from
\eqref{ado:eq:seed}. Expanding the denominator in the open disk is
justified by the integrable bound $|\kappa|/(1-|z|)$, and gives the
series \eqref{ado:eq:series}.

On a compact subset of $\Om$, $|1-zu|$ is bounded away from zero
uniformly in $u\in[0,1]$. The signed integral is consequently holomorphic
there and continues the same germ. If two finite signed measures have
these moments, their difference integrates every polynomial to zero.
Uniform polynomial approximation on $[0,1]$, followed by uniqueness of
a finite measure from its continuous-function integrals, proves equality.
\end{proof}

\begin{corollary}[Ordinary convergence on the nontrivial unit circle]
\label{ado:cor:boundaryseries}
For \eqref{ado:eq:domain}, $c_{\svec}(n)\to0$ and
\begin{equation}\label{ado:eq:coefficientvariation}
 \sum_{n\ge1}|c_{\svec}(n+1)-c_{\svec}(n)|
 \le\norm{\kappa_{\svec}}_1\le2^{d-1}.
\end{equation}
Thus the original series \eqref{ado:eq:series} converges at every
$|z|=1$, $z\ne1$, and agrees there with \eqref{ado:eq:signedtransform}.
\end{corollary}
\begin{proof}
Dominated convergence proves the coefficient limit. Also
\[
 |c(n+1)-c(n)|\le\int_0^1u^{n-1}(1-u)|\kappa(u)|\dd u.
\]
Summing the positive majorants gives \eqref{ado:eq:coefficientvariation}.
Summation by parts against the bounded partial sums of $z^n$ then proves
convergence when $|z|=1$, $z\ne1$. Abel's limiting theorem identifies its
sum with the radial limit of the disk series, which in turn equals the
continuous signed integral at that point.
\end{proof}
This conclusion concerns the parameter domain of this article. It is not
a claim that all of the manuscript's more general fractional-order
boundary series converge ordinarily.

\section{Exact oscillation and two interlacing laws}
\label{ado:sec:oscillation}
The central mechanism is a balance between moment orthogonality and
Rolle's theorem. Vanishing moments force sign changes; the two derivative
identities prevent extra zeros. This proof does not assume a
variation-diminishing theorem for fractional integration.

\begin{lemma}[Vanishing moments force sign changes]
\label{ado:lem:momentchanges}
Suppose a nonzero continuous real $f\in L^1(0,1)$ has finitely many
zeros and
\[
 \int_0^1u^j f(u)\dd u=0\qquad(0\le j\le m-1).
\]
Then $f$ has at least $m$ sign changes in $(0,1)$.
\end{lemma}
\begin{proof}
If its sign changes were $t_1,\ldots,t_k$ with $k<m$, the polynomial
$Q(u)=\prod_{j=1}^k(u-t_j)$ would have degree at most $m-1$ and
$Qf$ would have one fixed nonzero sign except at finitely many points.
Its integral would be nonzero, contradicting orthogonality. A zero at
which no sign change occurs is not included in $Q$.
\end{proof}

\begin{theorem}[Exact density zeros]
\label{ado:thm:oscillation}
For \eqref{ado:eq:domain}, $\kappa_{\svec}$ has exactly $d-1$
distinct simple zeros in $(0,1)$. Its sign alternates at those zeros,
and it is positive on the interval to the right of its largest zero.
\end{theorem}
\begin{proof}
At depth one, $f_b$ is strictly positive and has no zero. Suppose first
that a density $f$ at depth $d-1$ has exactly $d-2$ simple zeros.
By \eqref{ado:eq:opderivatives}, the derivative of $\Bop f$ has
exactly those zeros. Rolle's theorem therefore bounds the number of
distinct zeros of $\Bop f$ by $d-1$. This also rules out an infinite
set of zeros: any $d$ distinct zeros would already force too many
zeros of the derivative.

The first $d-1$ moments of the new density vanish, by
\eqref{ado:eq:opmoments}. Lemma~\ref{ado:lem:momentchanges} supplies
at least $d-1$ sign changes, so the upper bound is attained. There is
one derivative zero between each pair of consecutive new zeros, using
all $d-2$ derivative zeros. None remains at a new zero. Hence every
new zero is simple.

Now suppose $f$ is a depth-$d$ density with exactly $d-1$ simple
zeros, and put $g=\Aop f$. We have $g(1)=0$. If $g$ had $k$
interior zeros, Rolle's theorem between successive interior zeros and
between the largest one and the endpoint one would give at least $k$
zeros of $g'=-f/u$. Thus $k\le d-1$. Its first $d-1$ moments still
vanish, so the same lower bound gives exactly $d-1$ sign changes.
The derivative zeros are now all allocated to those successive
intervals, including the last interval ending at one. Consequently the
interior zeros of $g$ are simple. Repeating this argument for each
occurrence of $\Aop$ and $\Bop$ proves the zero count.

Let $P(u)$ be the monic polynomial whose roots are the resulting
$d-1$ zeros. The product $P\kappa_{\svec}$ has a fixed strict sign
off its roots. Orthogonality to degrees at most $d-2$ gives
\[
 \int_0^1P(u)\kappa_{\svec}(u)\dd u
 =\int_0^1u^{d-1}\kappa_{\svec}(u)\dd u=C_{\svec}>0.
\]
Thus $P\kappa_{\svec}$ is positive, and since $P$ is positive to
the right of its largest root, so is $\kappa_{\svec}$.
\end{proof}

\begin{corollary}[Index raising and depth raising interlace]
\label{ado:cor:interlacing}
If $f=\kappa_{\svec}$ has roots $r_1<\cdots<r_{d-1}$ and
$\Aop f$ has roots $q_1<\cdots<q_{d-1}$, then
\begin{equation}\label{ado:eq:Ainterlace}
 q_1<r_1<q_2<r_2<\cdots<q_{d-1}<r_{d-1}<1.
\end{equation}
Thus raising the first index by one moves every density zero strictly
toward zero. If a depth-$(d-1)$ density has roots
$t_1<\cdots<t_{d-2}$ and $\Bop f$ has roots $r_1<\cdots<r_{d-1}$,
then
\begin{equation}\label{ado:eq:Binterlace}
 r_1<t_1<r_2<\cdots<t_{d-2}<r_{d-1}.
\end{equation}
\end{corollary}
\begin{proof}
These are precisely the derivative-zero allocations in the preceding
proof. In the $\Aop$ case the last interval is $(q_{d-1},1)$; in the
$\Bop$ case all intervals are between consecutive interior zeros.
The number of derivative zeros equals the number of allocated intervals,
so the inequalities are strict and exhaustive.
\end{proof}

\subsection{Positive seeds, terminal mixtures, and a Lerch extension}
The preceding proof used only positivity, smoothness, integrability, and
the positive mass of the initial density. The explicit form
\eqref{ado:eq:seed} was needed to identify its moments as $n^{-b}$.

\begin{theorem}[Positive-seed extension]
\label{ado:thm:seedextension}
Let $f$ be a smooth strictly positive density in $L^1(0,1)$, of mass
$M_0>0$. Replace $f_b$ by $f$ in \eqref{ado:eq:kappaconstruction}.
All density-zero, interlacing, compensation, slit-plane nonvanishing,
and angular-count conclusions of this article remain valid. The norm
bound becomes $2^{d-1}M_0$, and the mass of the positive compensated
measure is
\[
 \frac{M_0}{d^{s_1}(d-1)^{s_2}\cdots2^{s_{d-1}}}.
\]
\end{theorem}
\begin{proof}
The operator identities still hold. After $d-1$ strict prefix steps,
the first $d-1$ moments vanish and the first nonzero moment is the
stated positive number. The proof of Theorem~\ref{ado:thm:oscillation}
therefore applies unchanged. The arguments in the next two sections
use only these moment and oscillation properties, and so establish the
remaining assertions as well.
\end{proof}
In particular, for fixed leading indices, every finite positive linear
combination
\begin{equation}\label{ado:eq:terminalmixture}
 \sum_{j=1}^N a_j L_{s_1,\ldots,s_{d-1},b_j}(z),
 \qquad a_j>0,\quad b_j>0,
\end{equation}
has the same nonvanishing and exact angular-count properties. Its seed
is $\sum a_jf_{b_j}$. This is a special convex-closure statement; sums
of arbitrary zero-free holomorphic functions need not be zero-free.

For a second application, take $\lambda>-1$ and
\[
 f_{b,\lambda}(u)=u^\lambda
                \frac{(-\log u)^{b-1}}{\Gamma(b)}.
\]
Its moments are $(n+\lambda)^{-b}$ and its mass is
$(1+\lambda)^{-b}$. The resulting strict shifted-terminal family is
\begin{equation}\label{ado:eq:lerchfamily}
 L^{(\lambda)}_{s_1,\ldots,s_{d-1};b}(z)
 =\sum_{n_1>\cdots>n_d\ge1}
 \frac{z^{n_1}}{n_1^{s_1}\cdots n_{d-1}^{s_{d-1}}(n_d+\lambda)^b}.
\end{equation}
It has the same minimal compensator degree $d-1$ and exactly $d-1$
angular zeros for $0<\rho\le1$. At depth one it is the usual
$z$-multiple of a Lerch series. No assertion is made here about shifts
in all denominators simultaneously.

\section{Canonical positive compensation and slit-plane nonvanishing}
\label{ado:sec:compensation}
Fix $\svec$ and abbreviate $\kappa=\kappa_{\svec}$,
$P=P_{\svec}$, $D=D_{\svec}$, and $C=C_{\svec}$. Empty products
at depth one are one. The preceding section proves that
\begin{equation}\label{ado:eq:positiveweight}
 P(u)\kappa(u)>0
 \quad\text{for all }u\in(0,1)\setminus\{r_1,\ldots,r_{d-1}\}.
\end{equation}
Since $P$ is bounded on $[0,1]$, the positive measure
$\dd\mu=P\kappa\dd u$ is finite.

\begin{lemma}[Polynomial cancellation]
\label{ado:lem:polycancel}
Suppose $m_1(\kappa)=\cdots=m_{d-1}(\kappa)=0$. For any real
polynomial $E$ of degree at most $d-1$, the removable quotient at zero
satisfies
\begin{equation}\label{ado:eq:generalcompensator}
 E(z)\frac{L(z)}{z^d}
 =\int_0^1\frac{u^{d-1}E(1/u)\kappa(u)}{1-zu}\dd u.
\end{equation}
Here $u^{d-1}E(1/u)$ is a polynomial, including at $u=0$.
\end{lemma}
\begin{proof}
First subtract the initial $d-1$ zero moments in
\eqref{ado:eq:signedtransform} to obtain
\[
 L(z)=z^d\int_0^1\frac{u^{d-1}\kappa(u)}{1-zu}\dd u.
\]
For $0\le j\le d-1$, multiplying this last transform by $z^j$
and moving the factor $z^j$ to the density changes it to
$u^{d-1-j}\kappa(u)$. The difference is a finite linear combination
of the moments of degrees $d-1-j,\ldots,d-2$, all zero. Applying
this monomial identity to every coefficient of $E$ proves
\eqref{ado:eq:generalcompensator} in the disk and hence throughout $\Om$.
\end{proof}

\begin{theorem}[Positive factorization and nonvanishing]
\label{ado:thm:factorization}
The measure $\mu$ has mass $C$, and
\begin{equation}\label{ado:eq:factorization}
 D(z)L(z)=z^d\int_0^1\frac{\dd\mu(u)}{1-zu}.
\end{equation}
Its transform $H$ has strictly positive imaginary part when
$\Imn z>0$, strictly negative imaginary part when $\Imn z<0$, and
$H(x)>0$ for every real $x<1$. The only zero of $L$ on $\Om$ is its
zero of multiplicity $d$ at zero.
\end{theorem}
\begin{proof}
We have $u^{d-1}D(1/u)=P(u)$, so the preceding lemma gives
\eqref{ado:eq:factorization}. Its value after division by $z^d$ at
zero is $\int P\kappa=C$, as already shown in the oscillation proof.
For nonreal $z$,
\[
 \Imn H(z)=(\Imn z)\int_0^1
              \frac{u\dd\mu(u)}{|1-zu|^2}.
\]
The integral is strictly positive because the measure has positive
density except at finitely many interior points. For real $x<1$,
the integrand defining $H(x)$ is positive. Thus $H$ has no zero on
$\Om$. All zeros of $D$ are $1/r_j>1$, outside $\Om$, while
$D(0)=1$ and $H(0)=C>0$. The factorization proves the claimed
multiplicity and the global nonvanishing assertion.
\end{proof}

\begin{theorem}[Minimality and uniqueness of the compensator]
\label{ado:thm:minimality}
Among real polynomials $E$ with $E(0)>0$ for which $E(z)L(z)/z^d$
is the transform of a finite positive measure on $[0,1]$, the smallest
possible degree is $d-1$. Every such polynomial of degree $d-1$ is
$E(0)D(z)$. In particular the normalization $E(0)=1$ singles out $D$.
\end{theorem}
\begin{proof}
Existence in degree $d-1$ follows from
Theorem~\ref{ado:thm:factorization}. Suppose that $\deg E\le d-1$.
Lemma~\ref{ado:lem:polycancel} already realizes the transform as the
finite signed measure with density
\[
 Q_E(u)\kappa(u),\qquad Q_E(u)=u^{d-1}E(1/u).
\]
Uniqueness of finite Hausdorff moments implies that any purported
positive representing measure must be this signed measure. Its density
must therefore be nonnegative almost everywhere. Since $\kappa$
changes sign at every $r_j$, continuity forces $Q_E(r_j)=0$.
As $r_j\ne0$, these are $d-1$ distinct roots $1/r_j$ of $E$.
Consequently $\deg E\ge d-1$. At equality those roots determine
$E$ up to its constant factor, which is fixed by $E(0)$.
\end{proof}
This is minimality of \emph{polynomial compensation for an ordinary
positive Stieltjes representation}. It is not minimal depth in an
algebra of numerical periods and not a minimal generalized Stieltjes
order. Those are different questions.

\subsection{Quantitative consequences}
For $x>0$, positivity and $0<u<1$ give
\begin{equation}\label{ado:eq:negativebounds}
 \frac{Cx^d}{(1+x)D(-x)}
 <(-1)^dL(-x)<\frac{Cx^d}{D(-x)}.
\end{equation}
Indeed $C/(1+x)<H(-x)<C$ and $D(-x)>0$. Without computing the
compensator roots one still has the coarser bounds
\begin{equation}\label{ado:eq:coarsebounds}
 C\left(\frac{x}{1+x}\right)^d
 <(-1)^dL(-x)<Cx^d,
\end{equation}
because $1\le D(-x)\le(1+x)^{d-1}$, with the relevant strictness
supplied by $H$ even at depth one.

Write $P(u)=\sum_{j=0}^{d-1}p_ju^j$ and
\begin{equation}\label{ado:eq:compensatedmoments}
 h_n=\int_0^1u^n\dd\mu(u)
     =\sum_{j=0}^{d-1}p_jc_{\svec}(n+j+1),\qquad n\ge0.
\end{equation}
For all $n,k\ge0$,
\begin{equation}\label{ado:eq:completepositivity}
 (-1)^k\Delta^kh_n
 =\int_0^1u^n(1-u)^k\dd\mu(u)>0,
 \qquad \Delta h_n=h_{n+1}-h_n.
\end{equation}
Every finite Hankel matrix $(h_{i+j})_{0\le i,j\le m}$ is positive
definite. The same is true of $(h_{i+j+1})$ and
$(h_{i+j}-h_{i+j+1})$. To see this, pair each matrix with a nonzero
real vector and integrate the square of its polynomial against,
respectively, $\mu$, $u\mu$, or $(1-u)\mu$. Each integral is strictly
positive. In particular,
\[
 h_nh_{n+2}-h_{n+1}^2>0.
\]
These are exact nonlinear positivity constraints on the original nested
harmonic coefficients and the compensator roots. They should not be
mistaken for rational-coefficient identities among special values.

\section{Exactly \texorpdfstring{$d-1$}{d-1} angular zeros on every upper semicircle}
\label{ado:sec:angular}
The signed transform gives, for $0<\rho\le1$ and $0<\theta<\pi$,
\begin{equation}\label{ado:eq:angulartransform}
 \Imn L(\rho e^{\ii\theta})
 =\rho\sin\theta\,J_\rho(\cos\theta),\qquad
 J_\rho(c)=\int_0^1
       \frac{\kappa(u)}{1-2\rho cu+\rho^2u^2}\dd u.
\end{equation}
A bound based only on simple sign-change arguments would leave open the
possibility of tangencies. We therefore prove a multiplicity-sensitive
upper bound before proving existence.

\subsection{An elementary confluent Cauchy-kernel bound}
\begin{lemma}[Angular variation bound with multiplicities]
\label{ado:lem:variation}
Suppose $\kappa\in L^1(0,1)$ is continuous, has exactly $m$ sign
changes at $r_1<\cdots<r_m$, and has a nonzero fixed sign on each
complementary interval. Then $J_\rho$ has at most $m$ zeros in
$(-1,1)$, counted with analytic multiplicity, for every $0<\rho\le1$.
\end{lemma}
\begin{proof}
The function
\begin{equation}\label{ado:eq:cauchymap}
 t(u)=\frac{1+\rho^2u^2}{2\rho u}
\end{equation}
is strictly decreasing on $(0,1)$, since
$t'(u)=(\rho^2u^2-1)/(2\rho u^2)<0$. Its values are at least one,
and it tends to infinity as $u\downarrow0$. Thus
\[
 J_\rho(c)=\int_0^1\frac{\kappa(u)}{2\rho u}
                         \frac1{t(u)-c}\dd u.
\]
The factor $\kappa(u)/u$ need not be integrable by itself; this causes
no difficulty because the displayed Cauchy factor is $O(u)$ at zero.
All the integrals below are interpreted with their complete kernels.

Suppose there are distinct zeros $c_1,\ldots,c_q\in(-1,1)$ with
positive multiplicities $\ell_1,\ldots,\ell_q$ summing to $m+1$.
If more zeros are available, retain any subcollection of total
multiplicity $m+1$. Form the proper rational function
\begin{equation}\label{ado:eq:variationrational}
 R(t)=\frac{\prod_{j=1}^m(t-t(r_j))}
             {\prod_{a=1}^q(t-c_a)^{\ell_a}}.
\end{equation}
Its partial fractions are linear combinations of
$(t-c_a)^{-h}$ for $1\le h\le\ell_a$. Differentiation of $J_\rho$
under the integral is valid locally for $c\in(-1,1)$; explicitly,
\[
 J_\rho^{(h-1)}(c)
 =(h-1)!\int_0^1\frac{\kappa(u)}{2\rho u}
                         \frac1{(t(u)-c)^h}\dd u.
\]
The selected zero multiplicities therefore imply
\begin{equation}\label{ado:eq:contradictoryintegral}
 \int_0^1\frac{\kappa(u)}{2\rho u}R(t(u))\dd u=0.
\end{equation}
This integral is absolutely convergent: near zero $R(t)=O(1/t)=O(u)$,
and near one its denominators stay positive and bounded away from zero.

But each factor $t(u)-t(r_j)$ changes sign exactly where $\kappa$
changes sign. The denominator of $R(t(u))$ is strictly positive, because
$t(u)\ge1>c_a$. Hence the integrand in
\eqref{ado:eq:contradictoryintegral} has one fixed nonzero sign except
at the $r_j$. Its integral cannot vanish. This contradiction proves the
bound and also excludes an identically zero $J_\rho$.
\end{proof}
The proof is a finite partial-fraction argument, including repeated poles.
No unproved strict-total-positivity assertion is being used.

\subsection{Positive compensation supplies the missing zeros}
For $0<\theta<\pi$ and $0\le u\le1$, put
\[
 \beta_u(\theta)=-\Arg(1-\rho u e^{\ii\theta}).
\]
For $0<u<1$ one has
\begin{equation}\label{ado:eq:betasector}
 0<\beta_u(\theta)<\beta_1(\theta)<\frac\pi2.
\end{equation}
The angle is strictly increasing in $u$, since differentiation with
respect to $u$ gives $\rho\sin\theta/|1-\rho u e^{\ii\theta}|^2>0$.
A positive mixture of the resolvents $(1-zu)^{-1}$ stays strictly inside
this sector. Consequently, with $H$ as in
\eqref{ado:eq:factorization},
\begin{equation}\label{ado:eq:Hsector}
 0<\Arg H(\rho e^{\ii\theta})<\beta_1(\theta).
\end{equation}
For example, the second inequality follows by rotating each resolvent
through $-\beta_1$ and observing that its imaginary part is strictly
negative for $0<u<1$. All rays lie in an angle of width less than
$\pi/2$, so no argument ambiguity arises.

\begin{theorem}[Complete angular count and signs]
\label{ado:thm:angular}
For \eqref{ado:eq:domain} and every $0<\rho\le1$,
$\Imn L(\rho e^{\ii\theta})$ has exactly $d-1$ simple zeros in
$0<\theta<\pi$. They have the signs and bounds stated in
Theorem~\ref{ado:thm:main}.
\end{theorem}
\begin{proof}
A continuous unwrapped argument of $L$ along the upper arc is
\begin{equation}\label{ado:eq:unwrappedphase}
 \Phi(\theta)=d\theta+
       \sum_{j=1}^{d-1}\beta_{r_j}(\theta)
       +\Arg H(\rho e^{\ii\theta}).
\end{equation}
The factorization proves this formula exactly; $\Phi$ is not reduced
modulo $2\pi$. The sector bounds imply
\begin{equation}\label{ado:eq:phasebounds}
 d\theta<\Phi(\theta)<d\bigl(\theta+\beta_1(\theta)\bigr)<d\pi.
\end{equation}
At the other end of the arc, $\Phi(\theta)\to d\pi$ as
$\theta\uparrow\pi$, since all resolvents approach positive real
values and $H(-\rho)>0$.

If $\rho<1$, then $\Phi(\theta)\to0$ as $\theta\downarrow0$.
The case $\rho=1$ requires more care: $L(1)$ may be singular, and
there is no reason to assign it argument zero. Instead, each fixed
$r_j<1$ gives $\beta_{r_j}(\theta)\to0$, while
\eqref{ado:eq:Hsector} gives
\[
 \limsup_{\theta\downarrow0}\Arg H(e^{\ii\theta})\le\pi/2.
\]
It follows that $0<\Phi(\theta)<\pi$ for sufficiently small positive
$\theta$. This weaker endpoint statement is all that is needed.

By continuity, the phase crosses each of $\pi,2\pi,\ldots,(d-1)\pi$.
These yield at least $d-1$ distinct zeros of the imaginary part. On the
other hand, Lemma~\ref{ado:lem:variation} with $m=d-1$, together
with \eqref{ado:eq:angulartransform}, gives at most $d-1$ zeros
counted with multiplicities. The change of variable $c=\cos\theta$
has nonzero derivative in the open arc, and the prefactor
$\rho\sin\theta$ is positive. Thus multiplicities are preserved.
Every zero is simple, and there is exactly one crossing of each listed
phase level. The levels must occur in their stated order, since a
reversal would force a repeated crossing of an earlier level.

The initial phase lies in $(0,\pi)$, so the imaginary part starts
positive. Simplicity forces alternating signs thereafter. At the
$k$th zero, the phase is $k\pi$ and the modulus is nonzero, so the real
value has sign $(-1)^k$.

For the location bounds, write $\alpha_k=k\pi/d$. At a zero,
\eqref{ado:eq:phasebounds} gives
\[
 \theta_k<\alpha_k<\theta_k+\beta_1(\theta_k).
\]
The function $h(\theta)=\theta+\beta_1(\theta)$ is strictly
increasing, because
\[
 h'(\theta)=\frac{1-\rho\cos\theta}
                  {1-2\rho\cos\theta+\rho^2}>0.
\]
For $\rho<1$, solving $h(\theta)=\alpha$ gives
$\theta=\alpha-\arcsin(\rho\sin\alpha)$. At $\rho=1$ the same
expression is the solution when $\alpha>\pi/2$; when
$\alpha\le\pi/2$ it is zero, and the strict lower bound follows
from $\theta_k>0$. This proves \eqref{ado:eq:mainangles}.
\end{proof}
The argument proves the needed crossings without asserting that
$\Phi$ is globally monotone. Such an assertion is unnecessary here.

\begin{corollary}[Analytic motion and no interior collisions]
\label{ado:cor:analyticbranches}
For fixed integer leading indices, each angular branch is real analytic
in $(\rho,b)\in(0,1)\times(0,\infty)$, and extends real analytically
across $\rho=1$ locally at each of its nontrivial upper-arc zeros.
The density roots $r_j$ are real analytic in $b>0$. No two of the
angular branches collide for $0<\rho\le1$.
\end{corollary}
\begin{proof}
On compact subsets of $\Ren b>0$, the seed $f_b$ is holomorphic as
an $L^1$-valued function; factors $|\log(-\log u)|^j$ arising from
parameter differentiation have an integrable common majorant on smaller
compact subsets. The bounded operators preserve this holomorphic
parameter dependence. The signed transform is locally holomorphic in
$z\in\Om$ and in $b$. Simplicity of its angular zeros permits the
real analytic implicit-function theorem. The unit-radius zeros remain
away from the cut, so the same local argument crosses $\rho=1$.
The density derivative formulas give joint real analyticity in $u,b$
in the open interval, and their simple roots obey the same theorem.
Distinctness follows from Theorem~\ref{ado:thm:angular}.
\end{proof}
The local extension across radius one does not imply that the global
count persists at all larger radii; Section~\ref{ado:sec:audit} gives
an explicit obstruction.

\subsection{Strict radial motion of every angular branch}
The compensation yields more than a zero count. It also turns a
monotonicity property of ordinary positive Stieltjes transforms into
strict motion of all the strict-depth angular zeros.

\begin{lemma}[Radial argument monotonicity for positive transforms]
\label{ado:lem:radialPick}
Let $\nu$ be a nonzero finite positive measure on $[0,1]$ and
$Q(z)=\int(1-uz)^{-1}\dd\nu(u)$. For $0<\theta<\pi$ and
$0<\rho\le1$,
\begin{equation}\label{ado:eq:radialPick}
 \frac{\partial}{\partial\rho}\Arg Q(\rho e^{\ii\theta})\ge0.
\end{equation}
The argument is chosen continuously from $Q(0)>0$.
\end{lemma}
\begin{proof}
First take a finite measure with positive weights at distinct points
$0<u_1<\cdots<u_n<1$. The rational function
$\sum_j w_j/(t-u_j)$ is strictly decreasing between consecutive poles,
with limits $+\infty$ and $-\infty$ at the two ends. Its numerator
therefore has exactly one root $v_j\in(u_j,u_{j+1})$ in each such
interval. Factoring the numerator gives
\[
 Q(z)=Q(0)\frac{\prod_{j=1}^{n-1}(1-v_jz)}
                      {\prod_{j=1}^n(1-u_jz)}.
\]
Consequently the radial derivative of its argument is
\[
 \sin\theta\left\{f(u_1)+
          \sum_{j=1}^{n-1}\bigl(f(u_{j+1})-f(v_j)\bigr)\right\},
 \qquad f(t)=\frac{t}{1-2\rho t\cos\theta+\rho^2t^2}.
\]
Every summand is nonnegative, because
\[
 f'(t)=\frac{1-\rho^2t^2}
            {(1-2\rho t\cos\theta+\rho^2t^2)^2}\ge0
 \qquad(0\le t\le1).
\]

Approximate an arbitrary finite positive $\nu$ weakly by such finite
measures, with their masses converging to $\nu([0,1])$. For each
compact subset of $\Om$, the integrands defining the transform and
its first complex derivative are bounded and continuous uniformly in
$u\in[0,1]$. The transforms and derivatives converge locally uniformly.
Also $Q(\rho e^{\ii\theta})\ne0$: it has positive imaginary part
unless the measure is supported at zero, in which case it is a positive
constant. Passing to the limit in
$\Imn(e^{\ii\theta}Q'(z)/Q(z))$ proves the nonnegative derivative.
The same argument works at $\rho=1$ because $\theta$ is strictly
between zero and $\pi$.
\end{proof}

\begin{theorem}[Strict inward angular motion]
\label{ado:thm:radialmotion}
For $d\ge2$ and \eqref{ado:eq:domain}, every angular branch satisfies
\begin{equation}\label{ado:eq:strictmotion}
 \frac{\dd\theta_k}{\dd\rho}<0,
 \qquad 0<\rho\le1,\quad 1\le k\le d-1.
\end{equation}
The derivative at one is the one-sided derivative, or equivalently the
derivative of the local continuation. The same result holds for the
positive-seed and shifted-terminal families of
Theorem~\ref{ado:thm:seedextension}.
\end{theorem}
\begin{proof}
At fixed $\theta$, differentiate the unwrapped phase
\eqref{ado:eq:unwrappedphase} with respect to the radius:
\[
 \Phi_\rho
 =\sum_{j=1}^{d-1}
      \frac{r_j\sin\theta}{1-2\rho r_j\cos\theta+\rho^2r_j^2}
   +\partial_\rho\Arg H(\rho e^{\ii\theta})>0.
\]
The first sum is strictly positive since $d\ge2$ and $r_j>0$; the
last term is nonnegative by Lemma~\ref{ado:lem:radialPick}.
At the unique simple crossing of the phase level $k\pi$, one has
$\Phi_\theta>0$. Indeed the phase starts below that level, ends
above it, and crosses it exactly once; simplicity excludes a zero
phase derivative there. Implicit differentiation of
$\Phi(\theta_k(\rho),\rho)=k\pi$ now gives
\[
 \theta_k'(\rho)=-\frac{\Phi_\rho}{\Phi_\theta}<0.
\]
All ingredients also hold for a general positive seed.
\end{proof}
For depth two, this does \emph{not} settle the manuscript's conjecture
about $\eta(\rho)=\cos\theta(\rho)/\rho$: the derivative of that
quotient has a second term, and its sign is not implied by
$\theta'(\rho)<0$. The two kinds of radial monotonicity must remain
separate in the integration ledger.

\section{Small-radius and large-outer-index expansions}
\label{ado:sec:asymptotics}
The global count removes an ambiguity from local expansions: the zeros
found near the first Fourier-mode zeros are the entire list of angular
zeros, not just selected branches.

\begin{theorem}[Two-term radial expansion]
\label{ado:thm:radialexpansion}
Fix an index satisfying \eqref{ado:eq:domain} with $d\ge2$, and put
\[
 p=\frac{c_{\svec}(d+1)}{C_{\svec}},\qquad
 q=\frac{c_{\svec}(d+2)}{C_{\svec}},\qquad
 \alpha_k=\frac{k\pi}{d}.
\]
The $k$th angular branch extends real analytically to $\rho=0$, and
\begin{align}
 \theta_k(\rho)={}&\alpha_k
 -\frac{p}{d}\sin\alpha_k\,\rho \notag\\
 &+\left(\frac{d+1}{d^2}p^2-\frac{2q}{d}\right)
       \sin\alpha_k\cos\alpha_k\,\rho^2
 +O_{\svec}(\rho^3).
 \label{ado:eq:radialexpansion}
\end{align}
In particular, every branch initially moves strictly toward smaller
angles as the radius increases from zero.
\end{theorem}
\begin{proof}
Divide the imaginary part of the disk series by $C_{\svec}\rho^d$.
The resulting analytic equation is
\[
 0=\sin(d\theta)+p\rho\sin((d+1)\theta)
                   +q\rho^2\sin((d+2)\theta)+O(\rho^3).
\]
At $(\rho,\theta)=(0,\alpha_k)$ its $\theta$ derivative is
$d(-1)^k\ne0$, so the analytic implicit-function theorem applies.
There are $d-1$ such branches; for small positive radius they exhaust
the global list from Theorem~\ref{ado:thm:angular}.

Write $\theta=\alpha_k+A\rho+B\rho^2+O(\rho^3)$ and divide by
$(-1)^k$. The coefficients of $\rho$ and $\rho^2$ are respectively
\[
 dA+p\sin\alpha_k,
 \qquad dB+(d+1)pA\cos\alpha_k+q\sin(2\alpha_k).
\]
Solving gives \eqref{ado:eq:radialexpansion}. The first derivative is
negative because $p>0$ and $0<\alpha_k<\pi$.
\end{proof}
Higher Taylor coefficients are obtained by the same finite coefficient
recursion. At each order the new unknown occurs with coefficient
$d(-1)^k$, while all remaining terms depend only on earlier
coefficients and finitely many moments. This is an exact algebraic
recurrence, not a proof of any global radial sign.

\subsection{Uniformity as the first index increases}
Fix $d\ge2$ and positive integer middle indices
$s_2,\ldots,s_{d-1}$. Let the first index be $a\in\bbN$ and the
last be $b>0$. Define
\[
 T_N(b)=\sum_{N\ge n_2>\cdots>n_d\ge1}
       \frac1{n_2^{s_2}\cdots n_{d-1}^{s_{d-1}}n_d^b},
 \qquad T_*=T_{d-1}(b)>0.
\]
The latter number is a single-chain value and is independent of $b$.
At depth two, $T_N(b)=H_N^{(b)}$ and $T_*=1$.

\begin{theorem}[First displacement, uniformly in the final order]
\label{ado:thm:largeouter}
Set $\lambda=d/(d+1)$ and $\mu=d/(d+2)$. As $a\to\infty$ through
positive integers, uniformly for $b>0$, $0\le\rho\le1$, and
$1\le k\le d-1$,
\begin{equation}\label{ado:eq:largeouter}
 \theta_k(\rho)
 =\alpha_k-\frac{T_d(b)}{dT_*}\sin\alpha_k\,
                      \rho\lambda^a
       +O\bigl(\rho^2\mu^a\bigr).
\end{equation}
At $\rho=0$, the branch is interpreted by its analytic extension.
The implied constant and the starting threshold for $a$ may depend on
the fixed depth and middle indices, but not on $b$, $\rho$, or $k$.
\end{theorem}
\begin{proof}
The normalized imaginary part is
\begin{equation}\label{ado:eq:outereq}
 \sin(d\theta)+\frac{T_d(b)}{T_*}\rho\lambda^a
                   \sin((d+1)\theta)+R(\theta),
\end{equation}
where
\[
 R(\theta)=\sum_{n\ge d+2}\frac{T_{n-1}(b)}{T_*}
        \rho^{n-d}\left(\frac dn\right)^a\sin(n\theta).
\]
All inner denominators are at least one, so
$0\le T_{n-1}(b)\le n^{d-1}$ uniformly for $b>0$. Choose the fixed
integer $a_0=d+3$. For $a\ge a_0$, termwise differentiation is
absolutely justified and
\begin{equation}\label{ado:eq:outertail}
 |R(\theta)|+|R'(\theta)|\le K\rho^2\mu^a,
 \quad
 K=\frac2{T_*}\sum_{n\ge d+2}n^d
                        \left(\frac{d+2}{n}\right)^{a_0}<\infty.
\end{equation}
Also $T_d(b)/T_*\le d^{d-1}/T_*$, independently of $b$.

On disjoint fixed small neighborhoods of the $\alpha_k$, the derivative
of the first term of \eqref{ado:eq:outereq} is bounded away from zero.
Its perturbation tends to zero uniformly by \eqref{ado:eq:outertail}.
Evaluating at $\alpha_k\pm K_1\rho\lambda^a$, for a sufficiently
large fixed $K_1$, gives opposite signs for all sufficiently large $a$.
This follows from the linear term $d(\theta-\alpha_k)$ and the
uniform bounds on the perturbations; $\mu/\lambda<1$. Thus each
neighborhood contains a unique root with displacement
$\delta=O(\rho\lambda^a)$. The global count identifies these roots
as all the branches in their established order.

Expand \eqref{ado:eq:outereq} at $\alpha_k$ and divide by $(-1)^k$.
The root equation becomes
\[
 0=d\delta+\frac{T_d(b)}{T_*}\rho\lambda^a\sin\alpha_k
     +O\bigl(\delta^2+\rho\lambda^a|\delta|
                         +\rho^2\mu^a\bigr).
\]
Because $\lambda^2<\mu$, the previously obtained displacement bound
makes the error $O(\rho^2\mu^a)$. Division by $d$ proves
\eqref{ado:eq:largeouter}. All constants used above are independent
of the final order.
\end{proof}
At depth two and unit radius this recovers the leading scale
$(2/3)^a$ already present in the manuscript's more detailed depth-two
expansion~\cite{repoSigned}. The added content is the simultaneous
all-depth statement, the exact branch indexing, and uniformity over the
entire range $b>0$ and $0\le\rho\le1$.

\section{Explicit logistic kernels and positive integral identities}
\label{ado:sec:identities}
The compensator is generally implicit: its roots are the zeros of a
signed density. The all-ones family makes every part explicit and gives
an infinite collection of exact positive integral identities.

\subsection{The all-ones density and its roots}
Let $\one_d=(1,\ldots,1)$ and define
\[
 \ell(u)=\log\frac{u}{1-u},\qquad 0<u<1.
\]
The familiar identity
\begin{equation}\label{ado:eq:allonesLi}
 L_{\one_d}(z)=\frac{(-\log(1-z))^d}{d!}
\end{equation}
follows recursively from
$(1-z)L'_{\one_d}(z)=L_{\one_{d-1}}(z)$ and the vanishing initial
value at zero. We use it as a classical identity with a short proof,
not as a new reduction of multiple polylogarithms.

\begin{theorem}[Explicit signed kernels and logistic roots]
\label{ado:thm:logistickernels}
For every $d\ge1$,
\begin{align}
 \kappa_{\one_d}(u)
 &=\frac{\Imn(\ell(u)+\ii\pi)^d}{\pi d!}\notag\\
 &=\sum_{j=0}^{\lfloor(d-1)/2\rfloor}
 \frac{(-1)^j\pi^{2j}\ell(u)^{d-1-2j}}
 {(2j+1)!(d-1-2j)!}.
 \label{ado:eq:logistickernels}
\end{align}
Its increasing list of roots is
\begin{equation}\label{ado:eq:logisticroots}
 r_j=\frac1{1+\exp(\pi\cot(j\pi/d))},
 \qquad 1\le j\le d-1.
\end{equation}
In particular, $r_j+r_{d-j}=1$.
\end{theorem}
\begin{proof}
Consider the generating function
\begin{equation}\label{ado:eq:kernelgenerator}
 G(t,u)=e^{t\ell(u)}\frac{\sin\pi t}{\pi t}
       =\sum_{d\ge1}K_d(u)t^{d-1}.
\end{equation}
For $|\Ren t|<1$, its integral at moment one is the beta integral
\[
 \int_0^1G(t,u)\dd u
 =\frac{\sin\pi t}{\pi t}\Gamma(1+t)\Gamma(1-t)=1.
\]
The reflection formula supplies the last equality, with removable values
understood at zero. Thus $\int K_1=1$ and $\int K_d=0$ for $d\ge2$.
Coefficient extraction is legitimate on a smaller complex neighborhood
of zero, since $u^{\Ren t}(1-u)^{-\Ren t}$ and its logarithmic
parameter derivatives have an integrable majorant.

Now $K_1=1$ and
\[
 \frac{\dd}{\dd u}G(t,u)=\frac{t}{u(1-u)}G(t,u),
\]
so $K_d'=K_{d-1}/(u(1-u))$. This is exactly the derivative of
$\Bop K_{d-1}$. Their difference is constant, and both have zero
integral for $d\ge2$, so $K_d=\Bop K_{d-1}$. Hence these are the
constructed all-ones kernels. Expanding the exponential and sine gives
\eqref{ado:eq:logistickernels}.

The imaginary part of $(\ell+\ii\pi)^d$ vanishes when
$\Arg(\ell+\ii\pi)$ is one of $j\pi/d$. This gives
$\ell=\pi\cot(j\pi/d)$. Reversing the order of these real values
and solving $u/(1-u)=e^\ell$ yields \eqref{ado:eq:logisticroots}.
The symmetry follows from $\cot(\pi-x)=-\cot x$.
\end{proof}

\begin{proposition}[A continuous Mellin identity and Stirling moments]
\label{ado:prop:mellin}
For real $x>0$ and complex $t$ satisfying $-x<\Ren t<1$,
\begin{equation}\label{ado:eq:mellin}
 \int_0^1u^{x-1}e^{t\ell(u)}\frac{\sin\pi t}{\pi t}\dd u
 =\frac{\Gamma(x+t)}{\Gamma(x+1)\Gamma(1+t)}.
\end{equation}
The values at removable singularities are obtained by continuation.
For every integer $n\ge1$,
\begin{equation}\label{ado:eq:stirlingmoments}
 \int_0^1u^{n-1}\kappa_{\one_d}(u)\dd u
 =\frac{\bracks{n}{d}}{n!},
\end{equation}
where $\bracks{n}{d}$ is the unsigned Stirling number of the first kind,
zero for $d>n$.
\end{proposition}
\begin{proof}
The integral equals
$\sin(\pi t)B(x+t,1-t)/(\pi t)$. The beta formula and Gamma
reflection simplify it to \eqref{ado:eq:mellin}. At integer $x=n$,
the right side is
\[
 \frac{(t+1)(t+2)\cdots(t+n-1)}{n!}.
\]
Its coefficient of $t^{d-1}$ is $\bracks{n}{d}/n!$, proving
\eqref{ado:eq:stirlingmoments} by \eqref{ado:eq:kernelgenerator}.
\end{proof}
These beta and Stirling identities provide an independent normalization
check on the general signed-moment construction.

\subsection{An explicit infinite family of compensated identities}
For the roots \eqref{ado:eq:logisticroots}, define
\[
 P_d(u)=\prod_{j=1}^{d-1}(u-r_j),\qquad
 D_d(z)=\prod_{j=1}^{d-1}(1-r_jz).
\]
The integrable function $P_d\kappa_{\one_d}$ is strictly positive
away from its finitely many zeros. Combining
\eqref{ado:eq:factorization} and \eqref{ado:eq:allonesLi} gives the
exact identity
\begin{equation}\label{ado:eq:generalidentity}
 \boxed{\quad
 \int_0^1\frac{P_d(u)\kappa_{\one_d}(u)}{1-zu}\dd u
 =\frac{D_d(z)}{d!}
       \left(\frac{-\log(1-z)}z\right)^d,
 \qquad z\in\Om.
 \quad}
\end{equation}
At $z=0$ the removable value is $1/d!$. At $z=-1$ this becomes
\begin{equation}\label{ado:eq:negativeoneidentity}
 \int_0^1\frac{P_d(u)\kappa_{\one_d}(u)}{1+u}\dd u
 =\frac{D_d(-1)}{d!}\log^d2.
\end{equation}
The content is not a new formula for $L_{\one_d}$; it is the explicit
positive density, its canonical compensation, and the resulting
positive-integrand identities. They remain valid at every depth without
a numerical relation search.

\subsection{Depth three}
Put
\[
 h_3=\frac\pi{\sqrt3},\qquad
 q_3=\frac1{2+2\cosh h_3}.
\]
Then
\[
 P_3(u)=u^2-u+q_3,\qquad
 D_3(z)=1-z+q_3z^2,\qquad
 \kappa_{\one_3}(u)=\frac{\ell(u)^2-\pi^2/3}{2}.
\]
Both factors $P_3$ and $\ell^2-\pi^2/3$ are negative between the
same two roots and positive outside them. Hence their product is
nonnegative throughout the interval. Formula
\eqref{ado:eq:negativeoneidentity} gives
\begin{equation}\label{ado:eq:depththreeidentity}
 \boxed{
 \int_0^1\frac{(u^2-u+q_3)(\ell(u)^2-\pi^2/3)}{1+u}\dd u
 =\frac{2+q_3}{3}\log^3 2.}
\end{equation}
The companion normalization is
\begin{equation}\label{ado:eq:depththreenormalization}
 \int_0^1 (u^2-u+q_3)(\ell(u)^2-\pi^2/3)\dd u=\frac13.
\end{equation}
These identities are proved statements, not high-precision conjectures.

\subsection{Depth four}
Set
\[
 q_4=\frac1{2+2\cosh\pi},\qquad
 P_4(u)=\left(u-\frac12\right)(u^2-u+q_4).
\]
Now
\[
 D_4(z)=\left(1-\frac z2\right)(1-z+q_4z^2),\qquad
 \kappa_{\one_4}(u)=\frac{\ell(u)^3-\pi^2\ell(u)}6.
\]
The roots are the logistic images of $-\pi,0,\pi$, so again the
product $P_4\kappa_{\one_4}$ is nonnegative. The corresponding pair
of identities is
\begin{align}
 \int_0^1\frac{P_4(u)(\ell(u)^3-\pi^2\ell(u))}{1+u}\dd u
 &=\frac{3(2+q_4)}8\log^4 2,
 \label{ado:eq:depthfouridentity}\\
 \int_0^1P_4(u)(\ell(u)^3-\pi^2\ell(u))\dd u&=\frac14.
 \label{ado:eq:depthfournormalization}
\end{align}
All endpoint singularities are logarithmic powers and are integrable.
Thus no principal-value convention is implicit in these formulas.

\subsection{An explicit density-edge asymptotic}
The smallest all-ones density root gives a simple large-depth benchmark:
\begin{equation}\label{ado:eq:edgeasymptotic}
 \log r_1=-d+\frac{\pi^2}{3d}+\frac{\pi^4}{45d^3}
             +O(d^{-5})+O(e^{-d}).
\end{equation}
Indeed $\log r_1=-\pi\cot(\pi/d)
-\log(1+e^{-\pi\cot(\pi/d)})$, and the Taylor expansion of $\cot$
at zero gives the displayed terms. This concerns a density root, not
an angular zero. Symmetry supplies $1-r_{d-1}=r_1$.

\section{Exact rational evaluation and finite certificates}
\label{ado:sec:computation}
The density zeros need not be computed in order to obtain numerical
certificates. The unsigned norm bound $\norm\kappa_1\le2^{d-1}$
already gives an explicit geometric remainder for a centered
resolvent expansion. The proof works with real final order; the delivered
exact implementation accepts positive integer indices so that all
required moments are rational.

\subsection{A centered-moment expansion}
Define
\begin{equation}\label{ado:eq:centeredmoments}
 B_k=\int_0^1(2u-1)^k\kappa(u)\dd u
 =\sum_{j=0}^k(-1)^{k-j}\binom{k}{j}2^j c_{\svec}(j+1).
\end{equation}
Then $|B_k|\le M$, where $M=2^{d-1}$. For $\Ren z<1$, put
\[
 w=\frac{z}{2-z},\qquad |w|<1.
\]
The inequality is equivalent to $|z|<|2-z|$. The identity
$1-zu=(1-z/2)(1-w(2u-1))$ gives
\begin{theorem}[Geometric exact-evaluation certificate]
\label{ado:thm:evaluator}
For every integer $N\ge1$,
\begin{align}
 V_N(z)&=\frac{z}{1-z/2}\sum_{k=0}^{N-1}B_kw^k,
 \label{ado:eq:centeredvalue}\\
 |L(z)-V_N(z)|
 &\le\left|\frac{z}{1-z/2}\right|
          \frac{M|w|^N}{1-|w|}.
 \label{ado:eq:centerederror}
\end{align}
At $z=\ii$ and even $N=2m$, a purely rational bound is
\begin{equation}\label{ado:eq:gaussianerror}
 |L(\ii)-V_{2m}(\ii)|<\frac{2^d}{5^m}.
\end{equation}
For $|z|\le1/2$, the uniform bound is
\begin{equation}\label{ado:eq:halferror}
 |L(z)-V_N(z)|\le\frac{2^{d-1}}{3^N}.
\end{equation}
\end{theorem}
\begin{proof}
The centered geometric series converges uniformly for $u\in[0,1]$,
and its tail is bounded by $|w|^N/(1-|w|)$. Integrating against
$|\kappa|$ proves \eqref{ado:eq:centerederror}. At $z=\ii$,
$|w|=1/\sqrt5$ and $|z/(1-z/2)|=2/\sqrt5$. The multiplicative
constant is therefore $2M/(\sqrt5-1)<2M=2^d$ after extracting
$5^{-m}$. On $|z|\le1/2$, one has $|w|\le1/3$ and
$|z/(1-z/2)|\le2/3$, proving \eqref{ado:eq:halferror}.
\end{proof}
For integer indices, \eqref{ado:eq:centeredmoments} and
\eqref{ado:eq:centeredvalue} use only rational arithmetic at rational
complex $z$. An enclosing complex disk gives enclosing intervals for
each coordinate by adding and subtracting its radius. The evaluator
never treats a small residual as a proof of an identity.

\subsection{Implementation choices and replay}
The file \texttt{code/exact\_evaluator.py} first obtains the moments
by dynamic strict prefix sums. It computes $B_k$ with a common integer
denominator and a repeated-difference transform, rather than a
floating-point alternating binomial sum. Horner evaluation then gives
the exact rational center. The generic branch can upper-bound square
roots by dyadic rationals; the Gaussian and half-disk cases use the
closed rational bounds above. Near $\Ren z=1$, the fixed norm-bound
precision in the generic branch can be insufficient, in which case the
implementation raises an error rather than returns an uncertified value.
It is not advertised as a complete optimal evaluator on the entire slit
plane.

From the package root, the principal replay is
\begin{verbatim}
python code/verify.py
\end{verbatim}
It reads frozen rational certificates and independently recomputes them.
The optional command
\begin{verbatim}
python code/verify.py --regenerate
\end{verbatim}
uses floating point to \emph{propose} angular brackets, then requires
rational arithmetic to certify both endpoint signs before saving them.
The proposal stage uses a strict-disk series and does not decide any
accepted sign. This discovery/replay separation is deliberate.

\subsection{Gaussian values}
The delivered Gaussian certificate file contains eight mixed-depth
examples, all evaluated with $N=240$ centered terms. Each complex-disk
radius is smaller than $10^{-81}$. The decimal entries below are
readable midpoint diagnostics; the certificates are the rational
endpoints in \texttt{data/gaussian\_enclosures.json}.
\begin{center}\small
\begin{tabular}{@{}lrr@{}}
\toprule
Index & $\Ren L_{\svec}(\ii)$, approximately
      & $\Imn L_{\svec}(\ii)$, approximately\\
\midrule
$(2,1)$ & $-0.174715661158$ & $-0.113648917600$\\
$(1,1,1)$ & $0.099953993665$ & $-0.033577147847$\\
$(2,1,2)$ & $0.032383901801$ & $-0.026348629072$\\
$(1,2,1,1)$ & $0.001158016810$ & $0.008343470492$\\
$(2,1,1,2)$ & $0.001825088603$ & $0.006510815484$\\
$(1,1,1,1,1)$ & $-0.003396506582$ & $-0.001886779708$\\
$(2,1,2,1,1)$ & $-0.000326599679$ & $-0.000010475684$\\
$(1,1,1,1,1,1)$ & $0.000443168800$ & $-0.000335617002$\\
\bottomrule
\end{tabular}
\end{center}
In addition, the eight all-ones Gaussian values at depths one through
eight are checked against \emph{independent exact} intervals for
\[
 L_{\one_d}(\ii)=\frac{(-\tfrac12\log2+\tfrac{\ii\pi}{4})^d}{d!}.
\]
The logarithm and $\pi$ intervals are derived from different rational
series, as detailed in Appendix~\ref{ado:app:controls}. These eight
normalization controls are distinct from the eight stored mixed-depth
examples above.

\subsection{Certified angular brackets}
On the semicircle of radius $1/2$, set
\begin{equation}\label{ado:eq:rationalcircle}
 z(t)=\frac{1-t^2}{2(1+t^2)}+\ii\frac{t}{1+t^2},
 \qquad \theta=2\arctan t,\quad t>0.
\end{equation}
Rational $t$ gives rational real and imaginary coordinates. Every saved
root bracket has $t$-width at most $10^{-12}$; the corresponding
angle-width is less than $2\cdot10^{-12}$. With $N=104$, the
error bound \eqref{ado:eq:halferror} proves the two endpoint signs by
strict rational comparisons. For illustration:
\begin{center}\small
\begin{tabular}{@{}lcrr@{}}
\toprule
Index & $k$ & $10^{12}t_{\rm left}$ & $10^{12}t_{\rm right}$\\
\midrule
$(2,1,2)$ & 1 & 472099547153 & 472099547154\\
$(2,1,2)$ & 2 & 1487963538987 & 1487963538988\\
$(1,2,1,1)$ & 1 & 314972569528 & 314972569529\\
$(1,2,1,1)$ & 2 & 806524527583 & 806524527584\\
$(1,2,1,1)$ & 3 & 2019771345242 & 2019771345243\\
$(1,1,1,1,1,1)$ & 3 & 774596669241 & 774596669242\\
\bottomrule
\end{tabular}
\end{center}
The file \texttt{data/angular\_brackets.json} contains 24 brackets in
eight complete index families of depths two through six, together with
all exact coordinate intervals. The existence of a zero in each bracket
follows from the opposite signs. The theorem, rather than a scan,
proves uniqueness in every bracket and excludes all additional angular
zeros. For the displayed middle branch at depth six, the all-ones
formula independently gives
$\theta_3(1/2)=\arccos(1/4)$ and $t=\sqrt{3/5}$.

\subsection{What the exact suite checks}
The fresh replay reports 5,713 finite check groups:
\begin{center}\small
\begin{tabularx}{\linewidth}{@{}Xr@{}}
\toprule
Check family & Count\\
\midrule
Strict coefficients against direct finite nested enumeration & 1080\\
Initial zero moments and leading coefficients & 120\\
Centered moments: independent binomial formula and norm bound & 2880\\
Polynomial compensation identities for arbitrary test roots & 1200\\
All-ones moments against unsigned Stirling recurrence & 240\\
Logistic-kernel derivatives through depth twelve & 11\\
Depth-three/four symbolic kernel formulas & 2\\
Exact generalized-kernel counterexample checks & 2\\
Rejected negative controls & 2\\
Rational atomic radial-argument derivative signs & 80\\
Independent exact all-ones Gaussian controls & 8\\
Replayed stored Gaussian enclosures & 8\\
Replayed angular endpoint signs & 48\\
Replayed root brackets and complete families & $24+8$\\
\midrule
Total & 5713\\
\bottomrule
\end{tabularx}
\end{center}
The polynomial compensation tests deliberately use arbitrary rational
roots: they test the cancellation identity, \emph{not} positivity for
those roots. Positivity requires the actual density zeros and is proved
analytically. Similarly, finite atomic derivative checks do not replace
the interlacing-and-approximation proof of
Lemma~\ref{ado:lem:radialPick}.

\section{Audit findings, corrections, and failed shortcuts}
\label{ado:sec:audit}
The source review was targeted, not a line-by-line audit of the entire
canonical manuscript. It covered the signed-kernel,
real-order, positive-double-kernel, subcritical, local-radius, and
research-status sections. No substantive counterexample to the reviewed
depth-two kernel and zero theorems was found. The following two
editorial corrections are concrete; the later items are scope safeguards
for this continuation rather than allegations that the manuscript makes
those mistakes.

\subsection{Two notation corrections in the transport proof}
In \path{chapters/05-real-positive-kernel.tex}, at the pinned snapshot,
the proof of the Hurwitz--polylogarithm transport formula says to
specialize the regularized Hurwitz identity ``with $q=n$'', and then
invokes boundedness of ``$h$''~\cite{repoPositive}. The displayed
Hurwitz formula uses $x$ as its real argument, and the regularized
kernel is named $q(t)$. The intended text is therefore
\[
 \text{``with $x=n$''},\qquad
 \text{``boundedness of $q$''}.
\]
The theorem and the coefficient calculation are unchanged. The file
\path{integration/notation_corrections.patch} contains precisely
these two replacements, at source lines 102 and 104. It is an optional
proposed patch; it has not been applied to the remote repository.

\subsection{Positive higher-order resolvents need not be zero-free}
The following exact example shows why the oscillation step in this
article cannot be replaced by a generic positivity slogan:
\begin{equation}\label{ado:eq:badpositive}
 R(z)=z^4\left(1+\frac1{(1-z)^4}\right)
     =z^4\int_{[0,1]}\frac{\dd(\delta_0+\delta_1)(u)}{(1-uz)^4}.
\end{equation}
The measure is positive, yet
\[
 z_0=1-e^{-\ii\pi/4}
     =1-\frac1{\sqrt2}+\frac{\ii}{\sqrt2}
\]
is a nonzero zero of $R$ and
$|z_0|^2=2-\sqrt2<1$. Indeed $(1-z_0)^4=-1$. Thus a positive
generalized Stieltjes kernel of order four does not imply even disk
nonvanishing. Our ordinary positive transform appears only
\emph{after} the sign-change polynomial has been identified and the
initial moments have been canceled.

\subsection{The unit-radius endpoint cannot be assigned an arbitrary phase}
At $\rho=1$, the point $z=1$ lies on the removed cut. A proof of
angular existence that simply sets $\Arg L(1)=0$ may silently assume
a finite endpoint value or a particular asymptotic. The proof in
Section~\ref{ado:sec:angular} instead uses the sector bound
$\Arg H<\beta_1$ to show that the unwrapped phase is below $\pi$
near the initial endpoint. This suffices even when $L$ diverges there.
No theorem about a limiting singular value is smuggled into the argument.

\subsection{The angular count does not extend automatically beyond radius one}
For the elementary depth-two example
$L_{1,1}(z)=\tfrac12\log^2(1-z)$, a zero of its imaginary part in
the upper half-plane requires
\[
 |1-z|=1.
\]
On $|z|=\rho$, this is $\cos\theta=\rho/2$. For $\rho>2$ there
is no such open-upper-arc zero at all. In our proof, the map
$t(u)=(1+\rho^2u^2)/(2\rho u)$ ceases to be monotone on $(0,1)$
when $\rho>1$, and the derivative used in the radial argument lemma
can also change sign. These are genuine losses of hypotheses.
Local continuation of a simple zero through $\rho=1$ must not be
confused with a theorem for all radii.

\subsection{Integer leading orders, normalized radii, and arithmetic depth}
The operator $\Aop^{s_j-1}$ was iterated an integer number of times.
Its elementary Rolle argument does not automatically apply to an
arbitrary fractional power. The manuscript itself proves that finite
signed moment representations can fail in part of the positive
real-order depth-two range~\cite{repoReal,repoSubcritical}. Therefore
one cannot state the present $L^1$ construction for all positive real
indices without a separate theorem or a different representation.

Likewise, strict decrease of $\theta_k(\rho)$ does not imply a chosen
sign for $\cos\theta_k(\rho)/\rho$. The original normalized-radius
conjecture and the fractional turning-point issues remain separate.
Finally, the degree $d-1$ in the compensation theorem is not a claim
about minimal multiple-polylogarithm depth at a numerical argument.
Formal relation modules, arithmetic period independence, and positive
transform compensation are different mathematical objects.

\subsection{The remaining arithmetic candidates retain their status}
The reviewed README and discovery chapter distinguish the proved $S_4$
identity, the unresolved $S_6$ and new $S_8$ candidates, and an older
$S_8$ vector that was rigorously rejected~\cite{repoStatus}. The
present analytic theorem does not furnish a relation certificate for
$S_6$ or $S_8$. Its Gaussian enclosures also do not establish arithmetic
independence or any conjectural equality merely by small residuals.
These status labels should remain unchanged when the package is
integrated.

\section{Further research questions and proposed directions}
\label{ado:sec:research}
The proofs above suggest a more structured programme than an unrestricted
search for small numerical relations. The following questions separate
analytic geometry, constructive certificates, and arithmetic reductions.

\subsection{Fractional leading indices}
Determine the largest real-index domain on which a finite signed moment
density exists and has exactly $d-1$ sign changes. The first focused
case is depth three with a real first index and integer second index.
The all-positive depth-two results show that absence of a finite signed
measure need not imply failure of zero-freeness. Consequently there are
two different problems: extending the present density theorem, and
extending the zero theorem by a compensated or distributional method.

For real leading indices at least one, a plausible testable target is
that the same $d-1$ oscillation statement persists. This is a proposed
conjectural direction, not a claim supported here by an exhaustive
fractional-parameter experiment. A proof would need a substitute for
the derivative-zero allocation used by $\Aop$.

\subsection{Normalized angular motion and its higher-depth analogue}
Theorem~\ref{ado:thm:radialmotion} settles strict decrease of the angles,
but not the manuscript's finer quotient
$\eta(\rho)=\cos\theta(\rho)/\rho$. At depth two, retain the
existing integer-outer-index normalized-radius conjecture exactly as
stated in the manuscript. At higher depth, a more natural displacement
is $(\alpha_k-\theta_k(\rho))/\rho$. Determine its monotonicity
regions, transition loci, and whether compensator-root interlacing
constrains them. The explicit coefficient in
\eqref{ado:eq:radialexpansion} provides an exact local diagnostic,
not a global conclusion.

\subsection{Motion under changes of an index or the terminal seed}
Raising the first index gives strict interlacing of density zeros.
Does it also order every angular zero in the expected direction toward
$\alpha_k$? Formula \eqref{ado:eq:largeouter} proves the limiting
behavior but does not compare every pair of finite consecutive outer
indices. A positive answer would require comparing both the compensator
and its positive transform, not only the locations of the compensator
roots. Analogous questions for an inner index, the terminal shift
$\lambda$, and positive terminal mixtures are also well posed.

\subsection{Sharper variation norms and faster certificates}
The bound $\norm\kappa_1\le2^{d-1}$ is uniform and easy to certify,
but it is not asserted optimal. Determine the best constant at fixed
depth, or at fixed depth and first index. The exact logistic kernels
supply a useful family in which the variation can be expressed through
one-dimensional integrals between known roots. A smaller certified
variation bound would directly tighten \eqref{ado:eq:centerederror}.
Another target is a full-slit evaluator with comparable exact
certificates and a conformal or Chebyshev expansion adapted to the cut.

\subsection{Constructive compensator roots and positive approximation}
Develop rational interval enclosures for all $r_j$ without assuming
closed forms. The proved interlacing laws can supply nested brackets
when an index or the depth changes. Once such enclosures are available,
combine the positive moment matrices of
\eqref{ado:eq:compensatedmoments} with quadrature or Pad\'e-type
approximants to produce two-sided bounds on real arguments. The
separation between an exact root theorem and a finite root-isolation
certificate should be retained throughout.

\subsection{Large-depth geometry}
For the all-ones family, \eqref{ado:eq:logisticroots} and
\eqref{ado:eq:edgeasymptotic} determine the density-edge scale.
For general index patterns, determine whether the edge roots still
have exponential scale in depth, and how the leading exponents depend
on the index profile. Compare this with the angular boundary layers
near $k=1$ and $k=d-1$. Uniform asymptotics with both $d$ and the
outer index growing would require estimates beyond the fixed-depth
remainder in Theorem~\ref{ado:thm:largeouter}.

\subsection{Radii beyond one and zero-count transitions}
The first target beyond the proved radius range is the explicit all-ones
family. There the problem reduces to the argument of
$-\log(1-\rho e^{\ii\theta})$. Classify the radii at which angular
zeros enter or leave through the cut endpoint, or form tangencies.
The elementary depth-two loss at $\rho=2$ shows that a universal
continuation without transitions is impossible. The higher-depth
problem may distinguish endpoint transitions from interior folds.

\subsection{Colored arguments and arithmetic identities}
For nontrivial inner colors, the real positive seed or the real signed
kernel can disappear. Identify color patterns admitting a real
character projection with controlled sign variation. This would be a
meaningful route toward analogous zero theorems for cyclotomic multiple
polylogarithms. Separately, pursue the manuscript's frozen $S_6$ and
$S_8$ candidates through exact functional-equation certificates. The
all-depth evaluator can supply checks, but a numerical enclosure
containing zero is not the missing arithmetic proof.

\subsection{A finite-to-analytic formalization path}
The first formal targets are the two operator moment identities,
finite polynomial cancellation, rational partial fractions with repeated
poles, and the exact evaluator's remainder arithmetic. The analytic
layer can then isolate Rolle allocation, Hausdorff-measure uniqueness,
positive-transform sectors, and the implicit-function argument.
None of the current code is represented as a Lean or other
proof-assistant formalization. A future formalization should record
exactly where classical compactness and analytic continuation enter,
rather than labeling the numerical replay as a substitute.

\appendix
\section{Independent rational logarithm and Gaussian controls}
\label{ado:app:controls}
The exact controls use only rational arithmetic and elementary remainder
bounds. They do not share the centered-moment evaluator's recurrence.
For $0<x<1$, let
\[
 A_N(x)=\sum_{j=0}^{N-1}\frac{(-1)^jx^{2j+1}}{2j+1}.
\]
The alternating-series estimate gives
\[
 \arctan x\in
 \begin{cases}
 [A_N(x),\ A_N(x)+x^{2N+1}/(2N+1)],&N\text{ even},\\
 [A_N(x)-x^{2N+1}/(2N+1),\ A_N(x)],&N\text{ odd}.
 \end{cases}
\]
The identity
\begin{equation}\label{ado:eq:machin}
 \pi=16\arctan(1/5)-4\arctan(1/239)
\end{equation}
is verified by tangent addition: if $a=\arctan(1/5)$, then
$\tan(4a)=120/119$, and
$\tan(4a-\arctan(1/239))=1$. Elementary bounds put the angle in
$(0,\pi/2)$, so it equals $\pi/4$ and not another tangent branch.
The implementation uses $N=120$ and $N=40$ for the two arctangent
series.

For the logarithm, set
\[
 B_N=2\sum_{j=0}^{N-1}\frac{(1/3)^{2j+1}}{2j+1}.
\]
Integrating the geometric series for $2/(1-x^2)$, or using the
power series of $2\operatorname{arctanh}x$, gives
\begin{equation}\label{ado:eq:log2interval}
 B_N<\log2<
 B_N+\frac{2(1/3)^{2N+1}}{(2N+1)(1-1/9)}.
\end{equation}
The code uses $N=160$. Both interval endpoints are exact rational
numbers. Cartesian complex interval multiplication then encloses
$(-\tfrac12\log2+\tfrac{\ii\pi}{4})^d/d!$.

For each depth $1\le d\le8$, the independently generated real and
imaginary intervals lie wholly inside the corresponding coordinate
intervals returned by the centered evaluator with 240 terms. This is
an exact containment check, not a decimal comparison. It is a strong
normalization and branch control, while the identity
\eqref{ado:eq:allonesLi} still rests on its differential-equation proof.

The exact evaluator also supplies a dyadic upper bound on $\sqrt{x}$
for a nonnegative rational $x=a/b$. With an integer precision $p$, set
\[
 q=\frac{\lfloor\sqrt{\lfloor 2^{2p}a/b\rfloor}\rfloor+1}{2^p}.
\]
Then $q^2>x$. This is computed by integer square root and exact
comparison. If its use in the generic geometric-tail formula fails to
produce an upper bound smaller than one for $|w|$, the code stops
rather than accepting a possibly divergent bound.

\section{Dependency map and integration instructions}
\label{ado:app:integration}
The logical dependencies are short enough to audit separately from the
implementation:
\begin{center}\small
\begin{tabularx}{\linewidth}{@{}p{0.34\linewidth}X@{}}
\toprule
Result & Dependencies\\
\midrule
Moment construction & Gamma integral; Fubini; the two operator identities.\\
Exact density zeros & Initial zero moments; Rolle's theorem; endpoint value of $\Aop f$.\\
Positive factorization & Exact density zeros; monic sign polynomial; polynomial cancellation.\\
Minimality & Uniqueness of finite Hausdorff moments; sign changes.\\
Global nonvanishing & Positive transform signs; compensator roots lie on the removed cut.\\
Angular upper bound & Strict monotonicity of $t(u)$; confluent partial fractions.\\
Angular lower bound & Positive-transform sector; continuous unwrapped phase.\\
Strict radial motion & Rational positive-transform interlacing; weak approximation; simple upward phase crossings.\\
Exact numerical bounds & Norm bound; centered geometric series; rational arithmetic.\\
\bottomrule
\end{tabularx}
\end{center}
In particular, the numerical checks are downstream of the analytic
proofs, and the angular upper and lower bounds have separate arguments.

A suggested report destination is
\begin{quote}\small\ttfamily
Analysis/Polylogarithms/docs/reports/all-depth-signed-oscillation/
\end{quote}
The integration synopsis in
\texttt{integration/all\_depth\_synopsis.tex} is a short standalone
section, using prefixed macros and labels. It summarizes the proved
scope and the proof chain; the full proofs remain in this article.
It should be placed after the existing signed-kernel and all-positive
depth-two developments, with the report cited as a separate contribution.
The optional notation patch is independent of this addition.

The manuscript's editorial ledger should distinguish four things:
ordinary proved theorems in this article, finite exact replay evidence,
proposed further questions, and previously open arithmetic candidates
whose status is unchanged. The word ``minimal'' should always retain
its qualifier ``polynomial compensation degree'' when discussing
Theorem~\ref{ado:thm:minimality}.

The package has no new dependency on a proof assistant, Wolfram Language,
or the full ProveIt checkout. Python's standard library suffices for the
exact evaluator. The replay additionally uses SymPy for the stated finite
symbolic checks; regenerating candidate angular brackets uses mpmath and
NumPy. The README records the versions actually used. A successful
rebuild may change PDF metadata; the package hash manifest describes the
delivered files, not every possible future rebuild.

\begin{thebibliography}{9}
\addcontentsline{toc}{section}{References}

\bibitem{repoSigned}
ProveIt repository, \emph{Polylogarithms manuscript},
Chapter 5, source \path{chapters/05-signed-kernels.tex}, especially
\texttt{signed:thm:kernel} and \texttt{signed:thm:unique-zero}.
Repository snapshot \texttt{\repoSHA}, October 2026.
Base directory:
\url{https://github.com/VladimirReshetnikov/ProveIt/tree/3447bc59b78d138a7f0d536b66ac6e5cc5d5e49f/Analysis/Polylogarithms/docs/manuscript}.

\bibitem{repoReal}
ProveIt repository, same manuscript and snapshot,
\path{chapters/05-real-orders.tex} and
\path{chapters/05-real-local-radius.tex}. Sources for the finite
signed-measure parameter domain and the local normalized-radius results.

\bibitem{repoSubcritical}
ProveIt repository, same manuscript and snapshot,
\path{chapters/05-subcritical-stieltjes.tex}. Sources for positive
consecutive differences, all-positive depth-two nonvanishing, and
subcritical radial geometry.

\bibitem{repoPositive}
ProveIt repository, same manuscript and snapshot,
\path{chapters/05-real-positive-kernel.tex}. Sources for the positive
double kernel and Hurwitz--polylogarithm transport. Blob SHA
\texttt{044d3825b90dce437ac9007733e6447cf9563098}.
The two notation replacements proposed here concern lines 102 and 104.

\bibitem{repoStatus}
ProveIt repository, same manuscript and snapshot,
\path{README.md} and \path{chapters/10-discovery.tex}. Sources for
proof-status distinctions, remaining arithmetic candidates, and the
surviving research programme. Consulted October 10, 2026.

\bibitem{LiuPego}
J.-G. Liu and R. L. Pego, \emph{On generating functions of Hausdorff
moment sequences}, Transactions of the American Mathematical Society
\textbf{368} (2016), 8499--8518, doi:10.1090/tran/6618.
The corrected author version arXiv:1401.8052v4, November 12, 2025,
including its corrigendum, was consulted:
\url{https://arxiv.org/html/1401.8052v4}.

\bibitem{LiStar}
Jiangtao Li, \emph{The derived set of multiple star polylogarithms},
arXiv:2609.03653v1, September 3, 2026.
\url{https://arxiv.org/html/2609.03653v1}.
This concerns weak/star sums; it is cited for comparison, not used to
replace the strict-sum proof.

\bibitem{DLMF}
NIST Digital Library of Mathematical Functions,
\S25.12, \emph{Polylogarithms}, and \S5.5, \emph{Functional Relations}
for the Gamma function.
\url{https://dlmf.nist.gov/25.12};
\url{https://dlmf.nist.gov/5.5}.
Consulted October 10, 2026.
\end{thebibliography}

\end{document}
